\documentclass[11pt]{article}

\usepackage[a4paper,margin=1.05in]{geometry}
\usepackage{amsmath,amssymb,amsthm}
\usepackage{mathtools}
\usepackage{stmaryrd}
\usepackage[protrusion=true,expansion=false]{microtype}
\usepackage{enumitem}
\usepackage[dvipsnames]{xcolor}
\usepackage{booktabs}
\usepackage{array}
\usepackage{listings}
\usepackage[hidelinks]{hyperref}

\hypersetup{colorlinks=true,linkcolor=BrickRed,citecolor=BrickRed,urlcolor=BrickRed}

\theoremstyle{plain}
\newtheorem{theorem}{Theorem}
\newtheorem{proposition}[theorem]{Proposition}
\newtheorem{lemma}[theorem]{Lemma}
\newtheorem{corollary}[theorem]{Corollary}
\newtheorem{conjecture}[theorem]{Conjecture}
\newtheorem{observation}[theorem]{Observation}
\theoremstyle{definition}
\newtheorem{definition}[theorem]{Definition}
\newtheorem{principle}[theorem]{Principle}
\newtheorem{obligation}{Obligation}
\theoremstyle{remark}
\newtheorem{remark}[theorem]{Remark}

\newcommand{\rhoc}{\textsf{rho}}
\newcommand{\V}{\mathbb{V}}
\newcommand{\Rnn}{\mathbb{R}_{\geq 0}}
\newcommand{\Cx}{\mathbb{C}}
\newcommand{\Bool}{\mathbb{B}}
\newcommand{\crisp}[1]{[\![#1]\!]}
\newcommand{\sem}[1]{\llbracket #1 \rrbracket}
\newcommand{\tensorv}{\otimes}
\newcommand{\sumv}{\oplus}
\newcommand{\dok}{\textsf{do}}
\newcommand{\ink}{\textsf{in}}
\newcommand{\fork}{\textsf{for}}
\newcommand{\where}{\textsf{where}}
\newcommand{\Cand}{\mathrm{Cand}}
\newcommand{\PA}{\mathrm{PA}}
\newcommand{\Tr}{\mathsf{Tr}}
\newcommand{\gag}{\mathcal{G}}
\newcommand{\Har}{\mathcal{H}}
\newcommand{\bisW}[1]{\approx_{#1}}
\newcommand{\treq}[1]{\simeq^{\mathrm{tr}}_{#1}}
\newcommand{\dia}[1]{\langle #1 \rangle}
\newcommand{\quo}[1]{@#1}

\lstdefinestyle{rholang}{
  basicstyle=\ttfamily\small,
  keywordstyle=\bfseries,
  morekeywords={for,new,in,where,contract,match,if,else,do},
  columns=fullflexible, keepspaces=true, showstringspaces=false,
  breaklines=true, frame=single, framesep=5pt, rulecolor=\color{black!35},
  xleftmargin=6pt, aboveskip=9pt, belowskip=9pt
}
\lstset{style=rholang}

\title{\bfseries A Gaggle of Transformers Under a Graded-Where Harness\\[4pt]
\large Continual learning by consolidating interventions, with the parameter
count held down by the causal structure the harness already owns}

\author{L.\,G.\ Meredith\\ \small F1R3FLY.io}
\date{F1R3FLY.io research note --- November 2026\\[2pt]\normalsize Draft 1}

\begin{document}
\maketitle

\begin{abstract}
The inhabiting-choice note established that a transformer trained on completed
runs is a rung-one object: a trace distribution has a unique autoregressive
factorisation, so branching structure is invisible to any learner that does
not act, and the agent loop around the model is where the acting happens. The
interventions-at-the-cut note then supplied the acting: a scoped $\dok$ on the
cut of the rho calculus with graded where-clauses, under which Pearl's calculus
embeds exactly and a refusal is a receipt that starves. This note puts the two
together as an engineering proposal. A \emph{harness} written in rho with
graded where-clauses and $\dok$ performs experiments on a population of
specimens; the experiments produce traces that carry their intervention, their
exogenous record, and their starvations; and the traces are consolidated into a
\emph{gaggle} of small transformers, each of which is itself a process whose
attention is a graded race in the log semiring. The claim defended is that a
gaggle \emph{indexed by interventions} is a rung-two object even though every
member of it is rung-one, because the family of interventional distributions
over a staged fragment determines the causal model, and that the harness,
being a rho term, is a witness in which the hedge is inherited --- which is
what moves the design off the saddle point the inhabiting-choice note located,
to the extent it moves at all. The parameter-count question gets a structural
answer rather than a compression trick: every gadget of the embedded causal
model needs a mechanism, not a world model, so a transformer in the gaggle
learns $P_i(v \mid \mathit{pa})$ over a parent set that is syntactic, with a
context window the size of that parent set and no obligation to represent the
exogenous noise the resolver already owns. Four further levers follow from the
companion notes: factorisation of the gaggle along namespaces the do-calculus
proves d-separated; pruning of attention heads by intervention on the
transformer's own races; consolidation of members that are \emph{interventionally}
bisimilar, a test the harness can afford because it owns their code; and
compilation of a member out of the gaggle altogether into an $\mathcal{F}_1$
clause when its learnt mechanism is a bounded-depth posynomial, at which point
the accumulated science lives in the clause language and the transformer was
scaffolding. The note closes with the algorithm as a rho sketch, the execution
story on f1r3node and Saph1r3, and seven obligations, of which the first ---
that the intervention-indexed family is rung two --- is the one the rest
depend on.
\end{abstract}

\tableofcontents

%%%%%%%%%%%%%%%%%%%%%%%%%%%%%%%%%%%%%%%%%%%%%%%%%%%%%%%%%%%%%%%%%%%%%%%%%%%%%%%
\section{The question}\label{sec:question}

The three companion notes draw a line, and the present note is an attempt to
build on the right side of it.

The first \cite{Inhabiting} observes that a system whose witnesses reduce
confluently can state branching phenomena but cannot inhabit them, and reads the
observation for artificial intelligence: a distribution over finite traces has a
unique autoregressive factorisation, so two systems that differ only in
\emph{when} a choice was made --- Milner's $\mathit{coin}.(\mathit{tea} +
\mathit{coffee})$ against $\mathit{coin}.\mathit{tea} +
\mathit{coin}.\mathit{coffee}$ --- induce the same trace distribution and are
indistinguishable to any learner whose only access is samples. The observable
that separates them is a refusal, and refusals do not appear in logs of completed
runs. A next-token predictor is therefore rung one of Pearl's ladder by an
information-theoretic limitation; the agent loop around it supplies rung two
from outside; and the arrangement --- confluent predictor, branching in the
harness --- is a saddle point, stable as long as per-query answers are what is
paid for and unstable the moment what is wanted is an accumulating, heritable
body of search structure.

The second \cite{Interventions} supplies the construct the harness needs in
order to intervene at all: $\dok(C)\ \ink\ P$, a scoped constraint on the
bindings of a race, placed at the cut because whether a cut will occur is
undecidable, lowered into the graded where-clause by the normal-form theorem for
weight maps, and exact on the staged fragment as Pearl's graph surgery. Under it
the vending-machine refusal is a receipt that starves, counterfactuals are
replay from a recorded pre-state, and the three rules of the do-calculus are
theorems about the embedded fragment.

The third family \cite{rhocomb3,matrix,matspec,saphire} is the route from a rho
term to a data-parallel device, through name-free combinators to an incidence
matrix of atoms against names, on which redex finding is a join, conflict
detection a second sparse product, and a machine step a maximal independent set
of the conflict graph. The route is strictly classical today --- no weights ---
and the conflict components it already computes are the contention sets at
which the graded semantics normalises.

The question this note asks is the one a working engineer asks after reading
the three: \emph{granted that a transformer is rung one and a harness in rho
with $\dok$ is rung two, what is the right way to lash several transformers
into such a harness so that the harness's experiments make the transformers
better, and so that the transformers stay small?} The emphasis on size is not
incidental. A transformer that must model the whole world is large because the
world is; a transformer that must model one mechanism of a causal model the
harness has already factored is small because the mechanism is. The size
question and the rung question turn out to have the same answer.

\subsection{The shape of the proposal}

\begin{itemize}[leftmargin=1.4em,itemsep=3pt]
\item A \emph{harness} $\Har$ is a rho term with graded where-clauses and
$\dok$, confined to a namespace it owns ($\mathcal{T}_4$ of
\cite{Interventions}), holding a budget. It performs \emph{experiments}: it
chooses a specimen, a pre-state, and a $\dok(C)$, runs the intervened term, and
records what happened.
\item A \emph{trace} is not a log of completed actions. It is a tuple
$(S, C, h, u, \star)$: the pre-state, the intervention, the causal history, the
exogenous record (which candidate fired at each site), and the starvation flag.
A trace carries its own intervention, which is what makes it rung-two data
rather than rung-one data (\S\ref{sec:data}).
\item A \emph{gaggle} $\gag = \{T_1, \dots, T_k\}$ is a population of small
transformers, each a process in the sense of the earlier transformer proposal
\cite{TransformerProposal}: attention is a graded receipt in the log semiring,
weights rest on channels, and the whole is a term the harness can quote,
intervene upon, and fund.
\item \emph{Consolidation} is the loop that takes traces to transformers:
routing a trace to the member that should learn from it, retraining that member
or spawning a new one, and then \emph{shrinking} the gaggle by merging members
the harness can prove interventionally equivalent and by compiling members out
into clauses when their learnt mechanism admits one.
\end{itemize}

\subsection{What the note claims, and what it concedes at once}

The claim is that this arrangement is a rung-two learner whose knowledge is
inherited, and that its parameter count is governed by the causal structure
the harness learns rather than by the size of the world.

The concession is that every transformer in the gaggle is, on its own, exactly
the rung-one object the inhabiting-choice note describes, and that the
branching lives in the harness. The note's answer to the saddle-point charge
is not that the charge is wrong; it is that the harness is a rho term, hence a
witness that can be a race, hence inheritable, and that the design's job is to
move as much of the accumulated science as possible from the transformers
(confluent, not inheritable as races) into the harness's clauses (races,
inheritable). \S\ref{sec:saddle} says exactly how far that goes.

%%%%%%%%%%%%%%%%%%%%%%%%%%%%%%%%%%%%%%%%%%%%%%%%%%%%%%%%%%%%%%%%%%%%%%%%%%%%%%%
\section{What the context fixes}\label{sec:context}

The note is self-contained in its definitions but not in its theorems; the
following are imported.

\begin{definition}[Graded receipt, contention, likelihood {\cite{GradedPLN}}]
A receipt $\fork(\vec y \leftarrow \vec x \;\where\; \psi)P$ is a join with a
clause $\psi$ valued in a commutative semiring $\V$, evaluated once per
candidate match. Two candidates contend when they share a message or a linear
receipt; a contention set is a connected component; resolution normalises within
it. Over $\Rnn$ the likelihood of a causal history $h = (\Gamma_i, \sigma_i)_i$
is $\prod_i w(\sigma_i)/\sum_{\tau \in \Gamma_i} w(\tau)$, and the parameter
gradient of its logarithm is a sum of local scores, one per firing site.
\end{definition}

\begin{definition}[The $\dok$ membrane {\cite{Interventions}}]
$\dok(y_1 \models \varphi_1 \,\&\, \cdots \,\&\, y_n \models \varphi_n)\ \ink\ P$
reads every receipt in scope binding $y_i$ with clause $\psi_r \tensorv
\crisp{\varphi_i}$, normalising after the restriction. On the flat fragment it
lowers exactly into the where-clause. On the staged fragment $\mathcal{T}_1$
every receipt is its own contention set, the causal order is a static DAG, and
$\dok$ on a gadget's ancilla binding is Pearl's graph surgery; positivity is the
absence of starvation; counterfactuals are replay with the exogenous record
held.
\end{definition}

\begin{proposition}[Observational invisibility {\cite[Prop.~4.1]{Inhabiting}}]
\label{prop:invisible}
Two labelled transition systems inducing the same trace distribution induce the
same autoregressive factorisation; a learner with access only to samples cannot
distinguish them.
\end{proposition}

\begin{proposition}[Interventions have one carrier {\cite[\S3]{Interventions}}]
\label{prop:carrier}
A site with a single candidate is intervention-transparent; the
observation/intervention gap lives only on races; over $\Rnn$ a context cannot
force a race it does not own.
\end{proposition}

\begin{observation}[The three requirements {\cite[Obs.~5.1]{Inhabiting}}]
\label{obs:rsi}
An improvement operator capable of recursion must (i) be defined on witnesses
that are races, (ii) preserve branching through evaluation, and (iii) fund
branches that lose at a rate that is itself subject to the operator.
\end{observation}

\begin{definition}[A transformer as a process {\cite{TransformerProposal}}]
\label{def:transformer}
A decoder-only transformer with $L$ layers and $H$ heads is the rho term in
which the residual stream at $(\ell, i)$ rests on a channel, keys and values at
$(\ell, h)$ rest as messages read by peek, and each query is a linear graded
receipt
\[
\fork(\,@j, @k, @v \lhd\!\!- \quo{[\mathrm{kv},\ell,h]}
 \;\where\; \crisp{j \leq i} \tensorv \tfrac{1}{\sqrt{d}}\, q\cdot k\,)\ P(v),
\]
read in the log semiring $(\mathbb{R}\cup\{-\infty\}, \mathrm{logaddexp}, +)$,
so that normalisation within the query's contention set is softmax and the
tropical limit is argmax. Under the path-sum reading, collapsed at the affine
boundary before the next nonlinearity, the receipt computes soft attention; under
the stochastic reading it computes stochastic hard attention with the same
expectation. The dense arithmetic is in the ground layer; the only race in the
term is attention.
\end{definition}

Definition~\ref{def:transformer} is what makes the rest of the note possible:
the transformer is a term, so the harness can quote it, run it under a $\dok$,
and read its exogenous record, and by Proposition~\ref{prop:carrier} the only
sites at which an intervention on it has content are its attention heads.

%%%%%%%%%%%%%%%%%%%%%%%%%%%%%%%%%%%%%%%%%%%%%%%%%%%%%%%%%%%%%%%%%%%%%%%%%%%%%%%
\section{The data: traces that carry their intervention}\label{sec:data}

The inhabiting-choice note remarks that ``logs that record failed attempts make
the refusal a trace symbol and the distinction learnable, but such logs are
intervention data in disguise, which is the point rather than a loophole.''
The present design takes the remark as its data schema.

\begin{definition}[Experimental trace]\label{def:trace}
An \emph{experimental trace} is a tuple $t = (S, C, h, u, \star)$ where $S$ is
a pre-state (a closed term, or its hash in the shard's state store), $C$ is a
constraint in $\mathcal{F}_0$ or $\mathcal{F}_1$, $h$ is the causal history of
$\dok(C)\ \ink\ S$ to its normal form or budget, $u$ is the exogenous record
(the fired class at each site), and $\star \in \{\mathrm{ok}, \mathrm{starved}\}$
records whether some restricted receipt starved. The \emph{observational
projection} of $t$ is the sequence of readouts along $h$; it is what a log of
completed runs would contain.
\end{definition}

The intervention $C$ is part of the datum, and so is the starvation flag. A
transformer trained on experimental traces is trained to predict $h$ \emph{given}
$(S, C)$, and that is a different object from a transformer trained on
observational projections.

\begin{proposition}[An intervention-indexed family is rung two]\label{prop:rungtwo}
Let $M$ be a positive causal Bayesian network over finite domains embedded as
the staged term $\sem{M}$, and let $D(C)$ denote the distribution of
observational projections of $\dok(C)\ \ink\ \sem{M}$. Then the family
$\{D(C) : C \in \mathcal{F}_0\}$ determines $M$ up to the equivalence of
causal Bayesian networks, whereas $D(\top)$ alone determines $M$ only up to
Markov equivalence. In particular the family separates the two vending machines
of \cite[\S4.6]{Inhabiting}, which $D(\top)$ does not.
\end{proposition}

\begin{proof}[Proof sketch]
By \cite[Thm.~5.3]{Interventions}, $D(C)$ for $C$ a point restriction on
gadget $i$ is the distribution of the submodel $M_{V_i = v}$. The joint
distributions of all single-variable hard interventions determine the DAG up to
the standard result that the interventional family identifies the graph when
every variable can be intervened on, and the conditional tables are read off the
intervened distributions directly. For the vending machines, $D(y = \mathit{tea})$
has a starved trace with positive probability for $Q$ and none for $P$, and the
starvation flag is part of the datum.
\end{proof}

The proposition is stated on the staged fragment because that is where the
companion note's theorems hold; Obligation~\ref{ob:rungtwo} asks for the
general statement. What it already says is enough to fix the design: \emph{the
harness must choose $C$, and must choose it well}. The quality of the gaggle is
bounded by the coverage of the interventions in its training set, and the
harness's job is experimental design, not data collection.

\begin{remark}[Which rung the transformer is on]
Each member of the gaggle remains a conditional trace predictor and is rung one
in the precise sense of Proposition~\ref{prop:invisible}: given $(S, C)$ it
predicts a trace distribution, and that distribution has a unique
factorisation. What is rung two is the \emph{family}, indexed by $C$, and the
index is supplied by the harness. This is the saddle point, named exactly, and
\S\ref{sec:saddle} is where it is addressed.
\end{remark}

%%%%%%%%%%%%%%%%%%%%%%%%%%%%%%%%%%%%%%%%%%%%%%%%%%%%%%%%%%%%%%%%%%%%%%%%%%%%%%%
\section{Where a transformer sits in a causal model}\label{sec:mechanism}

The embedding of \cite{Interventions} places, at each variable $V_i$, a gadget
whose clause is the conditional probability table $P_i(v \mid \mathit{pa})$ in
disjunctive normal form. On finite domains of modest size that clause is a
lookup. A transformer earns its place only where the lookup is unaffordable:
large or structured domains, mechanisms that are compositional in their inputs,
or variables whose values are themselves terms.

\begin{definition}[Transformer-backed gadget]\label{def:tgadget}
For a variable $V_i$ with parents $V_{p_1}, \dots, V_{p_m}$, a
\emph{transformer-backed gadget} replaces the enumerated clause by
\[
\fork(\,p_1 \leftarrow v^{(i)}_{p_1} \,\&\cdots\&\, p_m \leftarrow v^{(i)}_{p_m}
\,\&\, u \leftarrow w_i \;\where\; \crisp{\mathrm{dom}(*u)} \tensorv
\mathsf{T}_i(*p_1, \dots, *p_m)[*u]\,)\ \mathit{emit}_i(*u),
\]
where $\mathsf{T}_i$ is a member of the gaggle run as a pure ground expression
whose output is a vector of log-weights over $D_i$, indexed by the ancilla
value. The ancilla channel $w_i$ carries $D_i$ or, for large domains, is
sampled from the gaggle member's own vocabulary race.
\end{definition}

Three things about this placement decide the parameter question.

\begin{observation}[The context window is the parent set]\label{obs:context}
The clause at gadget $i$ reads $\mathit{pa}_i$ and nothing else, because the
parent set is syntactic \cite[\S1.2]{Interventions}: the names the pattern and
the clause read. A transformer backing gadget $i$ therefore has input length
$|\PA_i|$ (in tokens, the serialisation of the parent values), not the length of
the history. Attention's job of discovering which earlier positions matter has
been done by the harness's experiments and is recorded in the term as the
pattern; what remains for the transformer is the mechanism.
\end{observation}

\begin{observation}[The transformer does not model the noise]\label{obs:noise}
The exogenous noise $U_i$ is the resolver's draw at the ancilla race, and it is
the harness's, not the transformer's. A transformer trained on observational
projections must represent the correlations the noise induces across
variables --- confounding --- because it cannot tell them from mechanism. A
transformer backing a gadget predicts $P_i(v \mid \mathit{pa})$ and is
\emph{forbidden} by its placement from seeing anything else. The parameters that
an observational model spends on confounders are not spent.
\end{observation}

\begin{observation}[Only the races are learnable by intervention]\label{obs:races}
By Proposition~\ref{prop:carrier}, an intervention at a deterministic site is
either the identity or disables the site. The harness's experiments therefore
inform exactly the gadgets that are races --- the ancilla receipts --- and the
deterministic plumbing between them (fan-out, readout, the transformer's own
projections and MLPs) is intervention-transparent. This is not a limitation on
what the gaggle learns; it is a statement about where the gradient lands. The
score of \cite[Thm.~4.3]{GradedPLN} is a sum over firing sites, and a
deterministic site contributes zero.
\end{observation}

\begin{remark}[Activation patching is a $\dok$]
The interpretability community's causal tracing and activation patching ---
hold a run, replace one head's output with that from a corrupted run, measure the
change downstream --- is, in the vocabulary of this programme, replay under
$\dok(v \models \cdot)$ at an attention site with the exogenous record held,
which is \cite[Thm.~5.6]{Interventions} applied to the transformer's own term.
That the technique works on heads and is ill-defined on MLPs is
Proposition~\ref{prop:carrier}: heads are races, MLPs are not. The harness
therefore has, for free, an instrument for interrogating its own gaggle members
that the community built by hand, and \S\ref{sec:levers} uses it to prune.
\end{remark}

%%%%%%%%%%%%%%%%%%%%%%%%%%%%%%%%%%%%%%%%%%%%%%%%%%%%%%%%%%%%%%%%%%%%%%%%%%%%%%%
\section{The consolidation loop}\label{sec:loop}

\subsection{Overview}

The loop has five stages, each of which is a graded race in the harness term
and so is itself subject to the harness's budget and to intervention.

\begin{enumerate}[leftmargin=1.6em,itemsep=2pt]
\item \textbf{Propose.} Choose a specimen $S$ and a constraint $C$. The choice
is a graded race whose clause is the harness's current estimate of the
information the experiment would yield; \S\ref{ssec:propose}.
\item \textbf{Run.} Execute $\dok(C)\ \ink\ S$ under the budget; record the
experimental trace $t$ of Definition~\ref{def:trace}. On f1r3node this is one
deployment whose trace the validators already produce.
\item \textbf{Route.} Deliver $t$ to the gaggle member that should learn from
it. Routing is a graded receipt whose candidates are the members and whose
clause is each member's affinity for $t$; \S\ref{ssec:route}. If no member
accepts, the receipt starves, and starvation is the signal to spawn.
\item \textbf{Revise.} Retrain the chosen member on its accumulated traces, with
the PLN natural-gradient discipline of \cite{GradedPLN} so that confidence is
the metric; \S\ref{ssec:revise}.
\item \textbf{Shrink.} Merge members the harness can prove interventionally
bisimilar; prune heads by intervention; compile a member out into a clause when
its mechanism admits one; let members that no longer earn their keep starve;
\S\ref{sec:levers}.
\end{enumerate}

\subsection{Propose: experimental design as a graded race}\label{ssec:propose}

The harness holds, for each gadget $i$ of its current model, a graded clause
$\psi_i$ (enumerated or transformer-backed) and an evidence count $n_i$ in the
PLN sense. A candidate experiment is a pair $(S, C)$ with $C$ a point or set
restriction on some gadget's ancilla. Its value is placed in the clause of the
proposal receipt:

\begin{lstlisting}
for( @S <<- specimens & @C <<- candidates
     where  [ C in F0 and reachable(S, site(C)) ]
            * gain(S, C)                       // expected score, Eq. (1)
            * 1 / (1 + cost(S, C))             // phlo to run it
   ) { run!(S, C) }
\end{lstlisting}

The gain is the expected magnitude of the local score the trace would
contribute at the intervened site --- the quantity \cite[Cor.~5.7]{GradedPLN}
accumulates, available before the experiment is run because the candidate
weights the site would evaluate are the harness's own clause. Two remarks.
First, the proposal receipt is itself a race, so the harness hedges its
experiments: under the stochastic reading it does not always run the single
most informative one, and the fraction of budget spent on lower-gain
experiments is a clause constant the harness can tune --- requirement (iii) of
Observation~\ref{obs:rsi} applied to experimental design. Second, the
reachability guard is $\mathcal{F}_2$ in general and the vacuity question of
\cite[\S7.1]{Interventions}; under $\mathcal{T}_1$--$\mathcal{T}_3$ it is
decidable by inspection and is the positivity check.

\subsection{Route: the gaggle as a mixture of experts whose gate is a clause}
\label{ssec:route}

Members rest on a channel as quoted terms together with a \emph{signature}
$\mathrm{sig}_j$: a crisp formula over traces describing which traces the
member has been trained on --- in the simplest case, the set of sites its
gadgets back. Routing is

\begin{lstlisting}
for( @t <<- traces & @(j, sig, T) <<- gaggle
     where  [ sig(t) ] * affinity_j(t)
   ) { revise!(j, t) }
\end{lstlisting}

where $\mathrm{affinity}_j(t)$ is the likelihood member $j$ assigns to $t$'s
history under $t$'s intervention, evaluated by running $T_j$ as a pure
expression. Read in the log semiring this is a softmax gate over experts, which
is to say that the gaggle's routing \emph{is} attention over members, with the
same clause that attention within a member uses. Three consequences:

\begin{itemize}[leftmargin=1.4em,itemsep=2pt]
\item A trace no member's signature admits leaves the routing receipt with no
candidate of nonzero weight. The receipt starves. Starvation here is not death
but birth: a second receipt waits on the same channel with clause
$\crisp{\neg \bigvee_j \mathrm{sig}_j(t)}$ and spawns a fresh member with
signature generated from $t$'s site. The gaggle grows exactly when the world
presents a mechanism it has no expert for, and the \emph{refusal} of the gaggle
to absorb a trace is the growth signal --- positivity as starvation, read
constructively.
\item A trace several members admit is a race between them, and under the
stochastic reading more than one may be trained on it over time, which is the
hedge: the harness funds the losing expert at a rate set by the clause.
\item The signatures are crisp formulae in the term and are therefore visible
to the model checker, so ``is there a mechanism two members both claim?'' is a
spatial question the harness can ask of itself, and it is the trigger for the
merge test of \S\ref{ssec:merge}.
\end{itemize}

\subsection{Revise: retraining without losing plasticity}\label{ssec:revise}

A member retrained on its bucket faces the classical continual-learning
hazard: new traces overwrite old mechanisms. The companion paper measured the
same hazard in a different guise. Two learners circulating belief to one another
inflate their evidence counts without adding information, and the pathology is
\emph{loss of plasticity rather than instability} \cite[\S11]{GradedPLN}. The
natural-gradient result there --- revision is a Fisher-metric step with the
metric evaluated at the posterior count, hence implicit and unconditionally
stable --- is the discipline the gaggle should adopt, and it fixes two design
rules.

\begin{principle}[Conservation of evidence in the gaggle]\label{prin:evidence}
A member is never trained on a trace it generated. Traces are produced by
experiments on \emph{specimens}; a member's own outputs, when the harness runs it
as a specimen (\S\ref{ssec:merge}), are tagged as testimony and contribute to no
member's evidence count. Without this rule the gaggle inflates its own
confidence by re-reading itself, and the measured consequence is a gaggle that
can no longer learn.
\end{principle}

\begin{principle}[Confidence is the step size]
Each member carries, per gadget it backs, the PLN count $n_i$, and its learning
rate on traces at site $i$ is the Robbins--Monro schedule that count induces.
A member backing a well-established mechanism takes small steps; a fresh member
takes large ones. The count is data on a channel and the harness may intervene
on it --- resetting a count is the $\dok$ that restores plasticity.
\end{principle}

Both principles are the companion paper's results read for a population of
transformers rather than for a single clause, and neither requires anything a
transformer cannot do: the Fisher metric of a softmax output is standard, and
the per-site counts are bookkeeping.

%%%%%%%%%%%%%%%%%%%%%%%%%%%%%%%%%%%%%%%%%%%%%%%%%%%%%%%%%%%%%%%%%%%%%%%%%%%%%%%
\section{Keeping the parameters down}\label{sec:levers}

The brief asked specifically for the ability to keep the transformers small.
The answer has one structural part and four operational parts, and all five
are consequences of the companion notes rather than additions to them.

\subsection{The structural part: mechanisms, not worlds}

\begin{proposition}[Parameter bound on the staged fragment]\label{prop:bound}
Let $M$ be a causal Bayesian network over $n$ variables with domains
$D_1, \dots, D_n$ and parent sets $\PA_1, \dots, \PA_n$. A gaggle backing $M$
on the staged fragment needs, in total, parameters sufficient to represent
\[
\sum_{i=1}^{n} \;\Big|\textstyle\prod_{j \in \PA_i} D_j\Big| \cdot |D_i|
\]
table entries --- the size of the conditional probability tables --- and a
transformer is required at gadget $i$ only when that product is unaffordable as
a lookup. An observational model of the same system must represent the joint,
of size $\prod_i |D_i|$, less whatever factorisation it can infer from
observational data alone, which is the Markov-equivalence class and not $M$.
\end{proposition}

\begin{proof}
The first statement is Theorem~5.2 of \cite{Interventions}: the likelihood is
the truncated factorisation, and each factor is one gadget's clause. The second
is the standard gap between a joint and its CBN factorisation, with
Proposition~\ref{prop:rungtwo} supplying that interventions recover the
factorisation and observation does not.
\end{proof}

The proposition is almost a tautology, and that is its value: the parameter
count of the gaggle is the parameter count of the \emph{model}, in the sense in
which Pearl uses the word, and the harness's experiments are what establish the
model. The gaggle is small when the world is modular, and the harness finds out
whether it is.

\subsection{Lever 1: factorise along namespaces}

The composing-learners note \cite{Compose} established that a learner is a
namespace and that disjointness of namespaces is factorisation of the
distribution. The do-calculus gives the harness a test for disjointness it can
run: if $(Y \perp\!\!\!\perp Z \mid X)$ holds in the appropriate mutilated graph,
rule 1 of \cite[Cor.~5.4]{Interventions} says the mechanisms on $Y$ and $Z$ may
be learnt separately, and the harness should keep them in separate members.
The interventions note's first open problem --- d-separation as a crisp
behavioural formula over the staged term \cite[\S10]{Interventions} --- is
therefore load-bearing here:
when it is discharged, ``may these two members be kept apart?'' is a question
the model checker answers, and the gaggle's factorisation is maintained as a
checked invariant rather than a hope.

Two members on d-separated namespaces have, between them, the parameters of
two mechanisms. One member on their union has the parameters of a mechanism
over the product. The lever is to never pay for the product when the harness
has proved it need not.

\subsection{Lever 2: prune heads by intervention}

A member's attention heads are its only races (Definition~\ref{def:transformer}),
so they are the only sites at which the harness's $\dok$ has content on the
member itself. For each head $(\ell, h)$ of member $T_j$ the harness runs
\[
\dok(v \models \crisp{v = 0})\ \ink\ T_j \quad\text{at every query of head } (\ell,h)
\]
--- restrict the head's value race to the zero vector --- over a held-out set of
experimental traces, and compares the member's interventional likelihoods
before and after. A head whose removal leaves the likelihoods within tolerance
on every intervention in the held-out set is pruned. This is ordinary head
pruning; what the setting adds is that the comparison is over
\emph{interventional} distributions, so a head that matters only under
interventions the observational data never exhibited is kept, and one that
matters only for modelling confounders the member need not model is dropped.
The tropical limit of the same clause --- argmax attention --- is the
quantisation step of \cite{TransformerProposal} and composes with pruning.

\subsection{Lever 3: merge members that are interventionally bisimilar}
\label{ssec:merge}

The inhabiting-choice note's thesis, applied to the gaggle's own members, says
that trace equivalence between two members is the wrong test for merging them:
two members may assign identical likelihoods to every observational projection
and differ on what they do under an intervention. The right test is
bisimulation, and the harness can afford it because it owns the members' code
($\mathcal{T}_4$) and can run the bisimulation game on them.

\begin{definition}[Interventional bisimilarity of members]
Members $T_j$ and $T_k$ with the same signature are \emph{interventionally
bisimilar to tolerance $\epsilon$ on a constraint set $\mathcal{C}$} when for
every $C \in \mathcal{C}$ and every pre-state $S$ in the held-out set, the
distributions of experimental traces of $\dok(C)\ \ink\ T_j(S)$ and
$\dok(C)\ \ink\ T_k(S)$ are within $\epsilon$ in total variation, starvation
flags included.
\end{definition}

The test is the Larsen--Skou probabilistic bisimulation game
\cite{LarsenSkou91} with the harness as the tester, restricted to a finite
constraint set because the general question is undecidable, and the restriction
is honest: the members are bisimilar \emph{as far as the harness has been able
to experiment}, and the merge is an induction from that evidence with the
evidence count attached. Merging means retaining one member and deleting the
other, or distilling both into a fresh member of smaller size trained on the
union of their buckets under Principle~\ref{prin:evidence} (the distillation
targets are testimony, counted at zero).

The reason to prefer this test to the obvious one --- compare the members'
weights, or their outputs on a validation set --- is Proposition~\ref{prop:invisible}
once more. Output agreement is trace equivalence. Two members that agree on
every completed run and differ on a refusal are different mechanisms, and
merging them discards one.

\subsection{Lever 4: compile a member out into a clause}

The strongest lever is the one that removes a transformer entirely.

\begin{proposition}[Compile-out]\label{prop:compileout}
Let $T_j$ back gadget $i$ and suppose its learnt mechanism
$\mathit{pa} \mapsto P_i(\cdot \mid \mathit{pa})$ is representable, to tolerance
$\epsilon$ on the interventional held-out set, by a graded formula
$\psi \in \mathcal{F}_1$ of depth $d$ over the parent data. Then replacing the
transformer-backed gadget by the enumerated gadget with clause $\psi$ preserves
every interventional distribution of the model to tolerance $\epsilon$, costs
$O(2^d)$ per candidate at the site rather than a forward pass, and reduces the
member's parameter count to the constants of $\psi$.
\end{proposition}

\begin{proof}
The gadget's contention set is itself ($\mathcal{T}_1$), so its contribution to
every likelihood is its normalised clause value at the fired class, and two
clauses within $\epsilon$ on every parent value contribute within $\epsilon$
to every factor. The cost figure is \cite[Def.~7.2]{Interventions}.
\end{proof}

How to find $\psi$ is the research question, and the companion notes already
own half of it. An $\mathcal{F}_1$ graded formula over $\Rnn$ is a posynomial in
its constants \cite[Prop.~5.2]{GradedPLN}; fitting a posynomial of bounded
depth to a transformer's conditional outputs is a symbolic-regression problem
over a grammar the omnibus \cite{Omnibus} generates from the term language, and
the crisp subformulae are spatial tests on the parent data that the member's
attention pattern suggests. The proposal is to run compile-out as a standing
race in the harness: for each transformer-backed gadget, a candidate $\psi$ is
proposed from the member's attention structure, tested interventionally as in
\S\ref{ssec:merge}, and adopted if it passes, with the budget for the attempt a
clause constant.

What compile-out buys is more than parameters. A clause is a term; it is
quotable, inheritable by the reproduction note's crossover \cite{Reproduction},
checkable by the model checker, and \emph{a race}. A transformer's weights are
none of these. Every mechanism the gaggle manages to compile out has moved from
the confluent side of the line to the side where the accumulated science lives
as witnesses, which is the inhabiting-choice note's whole demand. The
transformers are scaffolding for mechanisms the harness has not yet understood
well enough to write down.

\subsection{Lever 5: mortality}

A member is funded from the harness's budget at a rate proportional to the
score its traces contribute. A member that backs a gadget no experiment has
touched within a window, or whose traces the merge test has shown another member
to cover, starves. This is the deadlock-by-starvation pruning of
\cite{ProofSearch} applied to the gaggle and needs no coordinator; the harness's
budget is conserved, and a member that does not earn tokens cannot pay for its
retraining. The funding rate for losing members --- those that lose routing
races --- is a clause constant and is itself in the harness's parameter set,
which is requirement (iii) of Observation~\ref{obs:rsi}.

\begin{table}[t]
\centering\small
\begin{tabular}{@{}>{\raggedright\arraybackslash}p{0.19\textwidth}>{\raggedright\arraybackslash}p{0.36\textwidth}>{\raggedright\arraybackslash}p{0.37\textwidth}@{}}
\toprule
\textbf{lever} & \textbf{what it removes} & \textbf{what licenses it}\\
\midrule
mechanisms, not worlds & parameters spent on the joint and on confounders & truncated factorisation; the parent set is syntactic\\
factorise & the product of two mechanisms when the sum suffices & do-calculus rule 1; d-separation as a checked formula\\
prune heads & heads that matter under no intervention & interventions have one carrier; heads are the races\\
merge & duplicate members & interventional bisimilarity, a test the harness can run\\
compile out & the transformer itself & $\mathcal{F}_1$ clauses are posynomials; exact on $\mathcal{T}_1$\\
mortality & members that earn nothing & conserved budget; starvation as pruning\\
\bottomrule
\end{tabular}
\caption{The parameter levers.}\label{tab:levers}
\end{table}

%%%%%%%%%%%%%%%%%%%%%%%%%%%%%%%%%%%%%%%%%%%%%%%%%%%%%%%%%%%%%%%%%%%%%%%%%%%%%%%
\section{The saddle point, revisited}\label{sec:saddle}

The inhabiting-choice note's charge against the loop architecture is precise:
the hedge lives in the harness, is collapsed at the boundary, and is handed to
no one as a race. The design above is a loop architecture and must answer the
charge on its own terms.

\paragraph{What is still on the saddle.} Each member is a confluent predictor.
Its parameters are improved by a gradient on outcomes, and the diversity its
next improvement would need is what the current improvement consumes, exactly
as \cite[\S5.4]{Inhabiting} says. Nothing here changes that about the
transformers.

\paragraph{What has moved.} Three things.

\begin{enumerate}[leftmargin=1.6em,itemsep=2pt]
\item The harness is a rho term. Its proposal race, its routing race, its
compile-out race are witnesses that are races, and its funding rates for the
losing branches of each are clause constants in the term. The harness
therefore meets the three requirements of Observation~\ref{obs:rsi} \emph{for
itself}: it can be quoted, inherited with its hedges uncollapsed, and improved
by an operator that acts on its clauses. The desert annual's germination
fraction is here the routing clause's constant for losing experts.
\item The gaggle is not a set of outcomes but a population of terms resting on
a channel, and the term that holds it carries the hedge: under the stochastic
reading of the routing race, several members are trained on overlapping
evidence and the population is passed on as a population. Best-of-$n$ collapses
at the boundary; a gaggle on a channel does not, because the consumer of the
gaggle is the harness, and the harness is a process runtime. This is the
``honest version'' of \cite[\S5.5]{Inhabiting}, built once as substrate.
\item Compile-out moves mechanisms across the line. A member that has been
compiled into a clause is no longer a confluent object; it is a race the
reproduction note can cross over, and the science it encodes is inherited as
science.
\end{enumerate}

\paragraph{The honest summary.} The design is a saddle point with a ramp off
it. The ramp is compile-out, and the design's success is measured by how much
of the gaggle, over time, it manages to replace with clauses. A gaggle that
never compiles anything out is the loop architecture with better data; a
gaggle that compiles everything out has become the mortal scientist of
\cite{Scientist}, and the transformers were the scaffolding it climbed.

%%%%%%%%%%%%%%%%%%%%%%%%%%%%%%%%%%%%%%%%%%%%%%%%%%%%%%%%%%%%%%%%%%%%%%%%%%%%%%%
\section{The harness, as a term}\label{sec:term}

The sketch below is in the concrete syntax of f1r3node's rholang with the
graded guard and $\dok$ as proposed in \cite{TransformerProposal,Interventions};
it is schematic in the ground-layer calls.

\begin{lstlisting}
new specimens, candidates, traces, gaggle, revise, run, budget in {

  // Propose: a graded race over experiments
  contract propose() = {
    for( @S <<- specimens & @C <<- candidates
         where [ C in F0 and positive(S, C) ] * gain(S, C) / (1 + cost(S, C))
       ) { run!(S, C) }
  } |

  // Run: one deployment under the do, record the experimental trace
  for( @S, @C <= run ) {
    do( C ) in { *S } |                 // the specimen, intervened
    for( @h, @u, @starved <- trace_of(S, C) ) {
      traces!((S, C, h, u, starved)) | propose!()
    }
  } |

  // Route: attention over members; starvation spawns
  for( @t <<- traces & @(j, sig, T) <<- gaggle
       where [ sig(t) ] * affinity(T, t)
     ) { revise!(j, t) } |
  for( @t <- traces where [ not(any_sig(t)) ] ) {
    new j in { gaggle!((*j, sig_from(t), fresh_member(t))) | revise!(*j, t) }
  } |

  // Revise: natural-gradient step with the per-site count as metric
  for( @j, @t <= revise ) {
    for( @(j, sig, T) <- gaggle & @n <- count(j, site(t)) ) {
      gaggle!((j, sig, natgrad_step(T, t, n))) | count(j, site(t))!(n + 1)
    }
  } |

  // Shrink: merge, prune, compile out, starve  (each a standing race)
  for( @(j, s, Tj) <<- gaggle & @(k, s, Tk) <<- gaggle
       where [ j != k ] * bisim_evidence(Tj, Tk) ) { merge!(j, k) } |
  for( @(j, s, T) <<- gaggle & @psi <<- clause_candidates(j)
       where interventional_fit(T, psi) ) { compile_out!(j, psi) } |
  for( @(j, s, T) <<- gaggle where [ earnings(j) < rent(j) ] ) { starve!(j) }
}
\end{lstlisting}

Everything the harness decides is a receipt with a clause, so everything it
decides can be intervened upon by a harness one level up, which is the
recursion that self-improvement needs. The sketch does not show the budget
flowing through each receipt; in the cost-accounted calculus every COMM is
priced and the seal tree of \cite{CostMonad} is where the flow is read off.

%%%%%%%%%%%%%%%%%%%%%%%%%%%%%%%%%%%%%%%%%%%%%%%%%%%%%%%%%%%%%%%%%%%%%%%%%%%%%%%
\section{Execution}\label{sec:execution}

\paragraph{On f1r3node.} The transformer proposal's Levels 0--2 are what the
gaggle needs: a ground tensor type, graded guards in the pure evaluator, and the
log-semiring resolver. The experimental trace of Definition~\ref{def:trace} is
what the order-then-execute design already records --- the executing validator
resolves the races and the winners are logged; the losers are recomputable from
the pre-state. The $\dok$ lowers at elaboration under $\mathcal{T}_2$. Nothing
in \S\ref{sec:loop} needs a construct the two companion proposals have not
already asked for, and the harness's own races are ordinary graded receipts.

\paragraph{On Saph1r3.} The Mat machine fires a maximal independent set of the
conflict graph per step, and the graded machine instead resolves each conflict
component by its algebra. The gaggle is a favourable workload: each member's
attention contention sets are star-shaped (one query, many keys) and small; the
routing race's contention set has one candidate per member; and the compile-out
race is embarrassingly parallel across members. The three items the earlier
proposal listed as missing from the route --- ground atoms in the Mat state,
peek as a primitive, and weights on the conflict graph --- are the same three
here, and the fourth, segmented reduction for the path-sum reading, is a
standard kernel. The route is nascent and the note does not pretend otherwise;
what it adds is that a gaggle of small members is a better first weighted
workload than one large one, because its contention sets are known statically
and its parallel step count is the depth of the harness's loop rather than the
width of the gaggle.

\paragraph{The budget.} \cite[Rem.~6.5]{GradedV4} asks that the denotation's
budget $\beta$, the risk ledger's $\omega$, the address-space bound, and the
inference-token budget be one budget doing several jobs. The gaggle adds the
parameter budget to the list: tokens buy parameters, parameters are paid for by
the score their traces earn, and the whole is one ledger.

%%%%%%%%%%%%%%%%%%%%%%%%%%%%%%%%%%%%%%%%%%%%%%%%%%%%%%%%%%%%%%%%%%%%%%%%%%%%%%%
\section{Obligations}\label{sec:obligations}

\begin{obligation}[The intervention-indexed family is rung two]\label{ob:rungtwo}
Proposition~\ref{prop:rungtwo} is stated on the staged fragment with
$\mathcal{F}_0$ constraints. The general statement should say, for a term off
$\mathcal{T}_1$, which constraint set $\mathcal{C}$ makes the family
$\{D(C) : C \in \mathcal{C}\}$ determine the term up to probabilistic
bisimulation, and should connect the answer to the branching cells of
\cite{AbbesBenveniste} as the interventions note's own Obligation 2 proposes.
Everything in \S\ref{sec:levers} rests on this.
\end{obligation}

\begin{obligation}[Merge soundness]
Show that replacing two interventionally $\epsilon$-bisimilar members on
constraint set $\mathcal{C}$ by one preserves the gaggle's interventional
distributions on $\mathcal{C}$ to a stated tolerance, and characterise the
tolerance's growth off $\mathcal{C}$.
\end{obligation}

\begin{obligation}[Compile-out search]
Give the grammar of $\mathcal{F}_1$ candidates generated from a member's
attention structure, and the acceptance test, as a graded race with a stated
budget; measure on the vending-machine benchmark and on the two laddered games
of \cite{Compose} how much of a gaggle compiles out and how fast.
\end{obligation}

\begin{obligation}[Evidence conservation under distillation]
Principle~\ref{prin:evidence} counts distillation targets at zero. The
testimony experiment of \cite[\S11]{GradedPLN} should be re-run with a
transformer-backed member in place of one clause to confirm that the
plasticity-loss pathology, and the remedy, transfer.
\end{obligation}

\begin{obligation}[The parameter bound off the finite fragment]
Proposition~\ref{prop:bound} is a CPT count. For continuous or term-valued
domains, where a transformer is needed, the bound should be restated in terms
of the member's capacity and the mechanism's description length in the clause
language, which is where compile-out's acceptance threshold comes from.
\end{obligation}

\begin{obligation}[Experimental design as a theorem]
The gain in \S\ref{ssec:propose} is the expected local score. The claim that
choosing experiments by it is information-optimal under the budget should be
stated against the optimal-experimental-design literature and against the
experimental-regimes note \cite{ExpRegimes}, whose interaction-only regimes
provably cannot certify universal negatives at polynomial cost --- a gaggle that
only ever intervenes without reflecting inherits that limit.
\end{obligation}

\begin{obligation}[The recursion]
\S\ref{sec:term} observes that every decision of the harness is a receipt a
harness one level up may intervene on. State the fixed point: a harness whose
specimens include its own quotation, with the budget as the unrolling depth, and
say what Observation~\ref{obs:rsi} demands of it that the one-level design does
not already supply.
\end{obligation}

%%%%%%%%%%%%%%%%%%%%%%%%%%%%%%%%%%%%%%%%%%%%%%%%%%%%%%%%%%%%%%%%%%%%%%%%%%%%%%%
\section{Conclusion}

A transformer cannot learn branching from traces, and a harness in rho with
graded where-clauses and $\dok$ can produce traces that carry their branching.
Lashing the two together is not a reconciliation of the inhabiting-choice note
with current practice; it is current practice done with the right data and an
honest account of where the race lives. The race lives in the harness, the
harness is a term, and the term is heritable --- that much is new. The members
of the gaggle remain confluent, and the design's ambition is to make them
temporary: to learn each mechanism with a small transformer whose context is a
parent set and whose job is one conditional, and then to compile it out into a
clause the moment the harness understands it well enough to write it down.

The parameter question answers itself along the way. A gaggle backing a causal
model is as large as the model's mechanisms and no larger; the harness's
experiments are what establish the factorisation that makes the mechanisms
small; the do-calculus tells the harness which members it may keep apart;
intervention on a member's own heads tells it which heads to keep; interventional
bisimulation tells it which members are one; and compile-out tells it when a
member can go. None of these is a compression technique applied from outside.
Each is the same observation --- that intervention sees what observation cannot
--- applied to the gaggle itself.

\section*{Acknowledgements}

The author of this note drafted it in dialogue with a transformer-derived
language model, which reviewed the companion notes and the Saph1r3 repository
and proposed the structure of \S\ref{sec:levers}. By the argument of
\cite[\S5]{Inhabiting} that is a rung-two interrogation of a rung-one model,
and the model's characterisation of its own members' situation in
\S\ref{sec:saddle} is reproduced with its agreement.

%%%%%%%%%%%%%%%%%%%%%%%%%%%%%%%%%%%%%%%%%%%%%%%%%%%%%%%%%%%%%%%%%%%%%%%%%%%%%%%
\begin{thebibliography}{99}\raggedright

\bibitem{Inhabiting} L.~G.~Meredith.
\newblock Inhabiting choice: confluent witnesses, the phenomena they cannot
carry, and one observation made three times.
\newblock F1R3FLY.io research note, November 2026
(\texttt{proof-dynamics-and-reasoning/inhabiting-choice}).

\bibitem{Interventions} L.~G.~Meredith.
\newblock Interventions at the cut: Pearl's \textsf{do} as a scoped guard in
the rho calculus with graded where-clauses.
\newblock F1R3FLY.io working note, draft 1, November 2026
(\texttt{fuzzyware/interventions-at-the-cut}).

\bibitem{GradedPLN} L.~G.~Meredith.
\newblock Graded where-clauses, II: the non-negative reals, and where the
mortal scientist gets gradient descent.
\newblock F1R3FLY.io, 2026 (\texttt{fuzzyware/graded-where-pln}).

\bibitem{GradedV4} L.~G.~Meredith.
\newblock Graded where-clauses, version 4.
\newblock F1R3FLY.io, 2026 (\texttt{fuzzyware/graded-where-v4}).

\bibitem{TransformerProposal} L.~G.~Meredith (with Claude).
\newblock A transformer in f1r3lang: an implementable model.
\newblock F1R3FLY.io proposal, September 2026.

\bibitem{rhocomb3} L.~G.~Meredith.
\newblock Name-free combinators for concurrency, draft 3.
\newblock F1R3FLY.io, 2026 (\texttt{rho-combinators/rho-combinators-draft3}).

\bibitem{matrix} L.~G.~Meredith.
\newblock The rho combinators as a matrix: a data-parallel representation of
the reduction graph, v1.2.
\newblock F1R3FLY.io, 2026 (\texttt{rho-combinators/rhocomb-matrix}).

\bibitem{matspec} L.~G.~Meredith.
\newblock F1R3Comb-Mat v0.5: a bulk-synchronous matrix machine.
\newblock F1R3FLY.io, 2026 (\texttt{rho-combinators/f1r3comb-mat-spec-v05}).

\bibitem{saphire} F1R3FLY.io.
\newblock Saph1r3: compiler from rho to GPU via rho combinators to matrix
representation.
\newblock \url{https://github.com/F1R3FLY-io/Saph1r3}, revision 205cdfb, 2026.

\bibitem{Compose} L.~G.~Meredith.
\newblock Composing learners: a learner is a namespace.
\newblock F1R3FLY.io, 2026 (\texttt{rho-life/compose-note}).

\bibitem{Scientist} L.~G.~Meredith.
\newblock The mortal scientist: learning as cost-accounted computation in a
reflective higher-order calculus.
\newblock F1R3FLY.io, 2026 (\texttt{rho-life/scientist-note}).

\bibitem{Reproduction} L.~G.~Meredith.
\newblock Reproduction: reflection-based genetics for mortal scientists.
\newblock F1R3FLY.io, 2026.

\bibitem{ProofSearch} L.~G.~Meredith.
\newblock Proof search as concurrent reaction.
\newblock F1R3FLY.io, 2026 (\texttt{proof-dynamics-and-reasoning/proof-search}).

\bibitem{ExpRegimes} L.~G.~Meredith.
\newblock Experimental regimes: laboratories for universal negatives, v3.
\newblock F1R3FLY.io, 2026 (\texttt{pnp-experiments/experimental-regimes-v3}).

\bibitem{Omnibus} L.~G.~Meredith.
\newblock Graph-structured lambda theories: a reference for the working
developer.
\newblock F1R3FLY.io, 2026.

\bibitem{CostMonad} L.~G.~Meredith.
\newblock Cost accounting as a monad.
\newblock F1R3FLY.io, 2026 (\texttt{cost-accounting-as-monad}).

\bibitem{LarsenSkou91} K.~G.~Larsen and A.~Skou.
\newblock Bisimulation through probabilistic testing.
\newblock \emph{Information and Computation} 94(1):1--28, 1991.

\bibitem{AbbesBenveniste} S.~Abbes and A.~Benveniste.
\newblock True-concurrency probabilistic models: branching cells and
distributed probabilities for event structures.
\newblock \emph{Information and Computation} 204(2):231--274, 2006.

\end{thebibliography}

\end{document}
